\documentclass[11pt]{article}
\usepackage[margin=1.1in]{geometry}
\usepackage{amsmath,amssymb,amsthm,mathtools,booktabs,array,tabularx,longtable}
\usepackage[dvipsnames]{xcolor}
\usepackage[colorlinks=true,linkcolor=MidnightBlue,citecolor=MidnightBlue,urlcolor=MidnightBlue]{hyperref}
\usepackage{fancyhdr}
\pagestyle{fancy}\fancyhf{}\renewcommand{\headrulewidth}{0pt}
\fancyfoot[R]{\thepage}
\fancyfoot[L]{\scriptsize\color{gray}Generated by the Transport Lab from approved parts of its store; not a deposited paper. Papers of record: \href{https://zenodo.org/communities/transport-papers}{Zenodo, Transport Papers}}
\fancypagestyle{plain}{\fancyhf{}\fancyfoot[R]{\thepage}\fancyfoot[L]{\scriptsize\color{gray}Generated by the Transport Lab from approved parts of its store; not a deposited paper. Papers of record: \href{https://zenodo.org/communities/transport-papers}{Zenodo, Transport Papers}}\renewcommand{\headrulewidth}{0pt}}
\newtheorem{theorem}{Theorem}[section]
\newtheorem{proposition}[theorem]{Proposition}
\newtheorem{lemma}[theorem]{Lemma}
\newtheorem{corollary}[theorem]{Corollary}
\newtheorem{conjecture}[theorem]{Conjecture}
\newtheorem{refutation}[theorem]{Refutation}
\newtheorem{prediction}[theorem]{Prediction}
\newtheorem{principle}[theorem]{Principle}
\newtheorem{observation}[theorem]{Observation}
\newtheorem{axiom}[theorem]{Axiom}
\theoremstyle{definition}
\newtheorem{definition}[theorem]{Definition}
\newtheorem{fact}[theorem]{Fact}
\newtheorem{assembled}[theorem]{Assembly}
\newtheorem{postulate}[theorem]{Postulate}
\newtheorem{premise}[theorem]{Premise}
\newtheorem{hypothesis}[theorem]{Hypothesis}
\newtheorem{remark}[theorem]{Remark}
\newtheorem{example}[theorem]{Example}
\theoremstyle{plain}
\setlength{\marginparwidth}{0.8in}
\newcommand{\changedmark}[2]{\marginpar{\raggedright\scriptsize\color{#1}#2}}
\providecommand{\Aa}{}\renewcommand{\Aa}{\mathcal{A}}
\providecommand{\Ad}{}\renewcommand{\Ad}{\widetilde{\mathrm{Ad}}}
\providecommand{\Aev}{}\renewcommand{\Aev}{\mathcal{A}^{+}}
\providecommand{\Aod}{}\renewcommand{\Aod}{\mathcal{A}^{-}}
\providecommand{\B}{}\renewcommand{\B}{\mathrm{B}}
\providecommand{\C}{}\renewcommand{\C}{\mathbb{C}}
\providecommand{\CAR}{}\renewcommand{\CAR}{\mathrm{CAR}}
\providecommand{\CS}{}\renewcommand{\CS}{{\mathbb{S}^{+}}}
\providecommand{\CSm}{}\renewcommand{\CSm}{{\mathbb{S}^{-}}}
\providecommand{\Cl}{}\renewcommand{\Cl}{\mathrm{Cl}}
\providecommand{\Con}{}\renewcommand{\Con}{\mathcal{C}}
\providecommand{\Dl}{}\renewcommand{\Dl}{\Delta}
\providecommand{\EN}{}\renewcommand{\EN}{E_{N}}
\providecommand{\End}{}\renewcommand{\End}{\mathrm{End}}
\providecommand{\F}{}\renewcommand{\F}{\mathbb{F}}
\providecommand{\G}{}\renewcommand{\G}{\Gamma}
\providecommand{\HH}{}\renewcommand{\HH}{\mathbb{H}}
\providecommand{\IS}{}\renewcommand{\IS}{I\sigma_3}
\providecommand{\J}{}\renewcommand{\J}{\mathbf{J}}
\providecommand{\LH}{}\renewcommand{\LH}{L_{H}}
\providecommand{\Mat}{}\renewcommand{\Mat}{\mathrm{M}}
\providecommand{\Mo}{}\renewcommand{\Mo}{M^{\circ}}
\providecommand{\Mp}{}\renewcommand{\Mp}{M^{+}}
\providecommand{\Msun}{}\renewcommand{\Msun}{M_{\odot}}
\providecommand{\Mt}{}\renewcommand{\Mt}{\widetilde{M}}
\providecommand{\Ncone}{}\renewcommand{\Ncone}{\mathcal{N}}
\providecommand{\Neff}{}\renewcommand{\Neff}{N_{\mathrm{eff}}}
\providecommand{\PCS}{}\renewcommand{\PCS}{\mathbb{P}{\mathbb{S}^{+}}}
\providecommand{\PT}{}\renewcommand{\PT}{\mathbb{PT}}
\providecommand{\Pin}{}\renewcommand{\Pin}{\mathrm{Pin}}
\providecommand{\Q}{}\renewcommand{\Q}{\mathbb{Q}}
\providecommand{\R}{}\renewcommand{\R}{\mathbb{R}}
\providecommand{\Ric}{}\renewcommand{\Ric}{\mathrm{Ric}}
\providecommand{\SBH}{}\renewcommand{\SBH}{S_{\mathrm{BH}}}
\providecommand{\Sc}{}\renewcommand{\Sc}{\mathbb{S}}
\providecommand{\Sch}{}\renewcommand{\Sch}{\mathrm{P}}
\providecommand{\Sig}{}\renewcommand{\Sig}{\Sigma}
\providecommand{\Som}{}\renewcommand{\Som}{S_\omega}
\providecommand{\Sp}{}\renewcommand{\Sp}{\mathbb{S}}
\providecommand{\Sph}{}\renewcommand{\Sph}{\mathrm{S}}
\providecommand{\Spi}{}\renewcommand{\Spi}{S_\pi}
\providecommand{\Spin}{}\renewcommand{\Spin}{\mathrm{Spin}}
\providecommand{\Sth}{}\renewcommand{\Sth}{S_{\mathrm{th}}}
\providecommand{\T}{}\renewcommand{\T}{\mathbb{T}}
\providecommand{\TH}{}\renewcommand{\TH}{T_{\mathrm{H}}}
\providecommand{\Tr}{}\renewcommand{\Tr}{\operatorname{Tr}}
\providecommand{\UHF}{}\renewcommand{\UHF}{\mathrm{UHF}}
\providecommand{\Z}{}\renewcommand{\Z}{\mathbb{Z}}
\providecommand{\Zc}{}\renewcommand{\Zc}{\mathrm{Z}}
\providecommand{\cmg}{}\renewcommand{\cmg}{\mathrm{cm^{2}\,g^{-1}}}
\providecommand{\dS}{}\renewcommand{\dS}{\mathrm{dS}}
\providecommand{\dsc}{}\renewcommand{\dsc}{\mathrm{dsc}}
\providecommand{\dtPG}{}\renewcommand{\dtPG}{dt}
\providecommand{\dtS}{}\renewcommand{\dtS}{dt_{\mathrm{S}}}
\providecommand{\ellv}{}\renewcommand{\ellv}{\ell}
\providecommand{\er}{}\renewcommand{\er}{e_r}
\providecommand{\et}{}\renewcommand{\et}{e_t}
\providecommand{\ev}{}\renewcommand{\ev}{\mathrm{ev}}
\providecommand{\gl}{}\renewcommand{\gl}{\mathfrak{gl}}
\providecommand{\hatV}{}\renewcommand{\hatV}{\widehat{V}}
\providecommand{\hb}{}\renewcommand{\hb}{\bar{h}}
\providecommand{\hf}{}\renewcommand{\hf}{\hat{f}}
\providecommand{\id}{}\renewcommand{\id}{\mathrm{id}}
\providecommand{\kpc}{}\renewcommand{\kpc}{\mathrm{kpc}}
\providecommand{\lP}{}\renewcommand{\lP}{\ell_{P}}
\providecommand{\mD}{}\renewcommand{\mD}{\mathbf{m}_{D}}
\providecommand{\mb}{}\renewcommand{\mb}{\bar m}
\providecommand{\mnu}{}\renewcommand{\mnu}{\Sigma m_{\nu}}
\providecommand{\nInf}{}\renewcommand{\nInf}{n_\infty}
\providecommand{\nO}{}\renewcommand{\nO}{n_0}
\providecommand{\ninf}{}\renewcommand{\ninf}{n_{\infty}}
\providecommand{\nz}{}\renewcommand{\nz}{n_{0}}
\providecommand{\proj}{}\renewcommand{\proj}{\operatorname{proj}}
\providecommand{\pv}{}\renewcommand{\pv}{\pi}
\providecommand{\rad}{}\renewcommand{\rad}{\mathrm{rad}}
\providecommand{\sgn}{}\renewcommand{\sgn}{\operatorname{sgn}}
\providecommand{\shA}{}\renewcommand{\shA}{\mathcal{A}}
\providecommand{\shC}{}\renewcommand{\shC}{\mathcal{C}}
\providecommand{\shE}{}\renewcommand{\shE}{\mathcal{E}}
\providecommand{\shJ}{}\renewcommand{\shJ}{\mathcal{J}}
\providecommand{\shK}{}\renewcommand{\shK}{\mathcal{K}}
\providecommand{\shL}{}\renewcommand{\shL}{L}
\providecommand{\shW}{}\renewcommand{\shW}{\mathcal{W}}
\providecommand{\sll}{}\renewcommand{\sll}{\mathfrak{sl}}
\providecommand{\so}{}\renewcommand{\so}{\mathfrak{so}}
\providecommand{\su}{}\renewcommand{\su}{\mathfrak{su}}
\providecommand{\supp}{}\renewcommand{\supp}{\operatorname{supp}}
\providecommand{\tr}{}\renewcommand{\tr}{\operatorname{tr}}
\providecommand{\uu}{}\renewcommand{\uu}{\mathfrak{u}}
\providecommand{\br}{}\renewcommand{\br}[1]{\langle #1\rangle}
\providecommand{\bra}{}\renewcommand{\bra}[1]{\langle #1|}
\providecommand{\file}{}\renewcommand{\file}[1]{\texttt{#1}}
\providecommand{\hp}{}\renewcommand{\hp}[2]{( #1 \,|\, #2 )}
\providecommand{\ip}{}\renewcommand{\ip}[2]{\langle #1, #2\rangle}
\providecommand{\ket}{}\renewcommand{\ket}[1]{|#1\rangle}
\providecommand{\rel}{}\renewcommand{\rel}[2]{D\bigl(#1 \,\|\, #2\bigr)}
\providecommand{\rev}{}\renewcommand{\rev}[1]{\widetilde{#1}}
\providecommand{\slashed}{}\renewcommand{\slashed}[1]{{#1}\mkern-10mu/}
\title{\bfseries Dark energy and the frozen point}
\author{Leon Butler\\ \normalsize\href{mailto:lbutler.math@gmail.com}{\texttt{lbutler.math@gmail.com}}\\ \normalsize ORCID: \href{https://orcid.org/0009-0006-4449-5617}{0009-0006-4449-5617}}
\date{Generated by the Transport Lab, 2 November 2026}
\begin{document}
\maketitle
\begin{abstract}
\normalsize
The conformal completion contains one distinguished fixed direction, the frozen point, and its length is set by the cosmological constant. This trail follows what that identification buys.

The frozen point supplies the extra vector that a known polynomial form of gravity needs, and its wobbles turn out to carry no physics. Where space is empty apart from the cosmological constant, its length cannot change, so dark energy behaves exactly like a cosmological constant.

On one stated condition this becomes a sharp prediction that dark energy never evolves, which a confirmed measurement of evolution would refute. The trail ends with what the framework does not do: it leaves the value of the cosmological constant, and even its sign, as an input.
\end{abstract}
\newpage
\noindent \textbf{Authorship and disclosure. }The founding idea of the series is the author's: that division by zero in a Clifford algebra should be read not as a prohibition but as a transport into a higher-dimensional algebra, and that the transported division need not be accompanied by an inverse operation. The direction of investigation and the critical review of the results are likewise the author's. The mathematical development and the drafting of the text were carried out in extended collaboration with Claude Fable 5 (Anthropic), an AI system.

\noindent\fbox{\parbox{\dimexpr\linewidth-2\fboxsep-2\fboxrule}{\textbf{This is not a deposited paper.} It was generated by the Transport Lab on 2 November 2026 and assembled from approved parts of its store. It follows results across several papers of the series; each is stated with its proof and credited to the paper that first proved it. The official versions of those papers are on Zenodo, in the Transport Papers community, \url{https://zenodo.org/communities/transport-papers}.}}
\newpage

\section*{Introduction}
The walk first builds the geometry of the frozen point, then reads off its prediction for dark energy, and finally asks what kind of number the cosmological constant is. Results a step rests on but that are not on the trail are listed at the end.

\section{What we assume}

\begin{postulate}[(FS), the field-space premise of the gravitational action: the varied gravitational connection, a $\Cl(2,4)$ gauge field with nondegenerate soldering, admits a parallel tractor $I$ with $\ip{X}{I} = 1$ (the physical scale), a holonomy reduction to the stabiliser of $I$]\label{c:synthesis.field-space-premise-parallel-frozen-point}
The series borrows a rule, an action, from which the equations of gravity follow, and that rule gives Einstein's equations only if one decides in advance which shapes of the gravitational field are allowed to vary: those that keep one special fixed direction, the frozen point, unchanged everywhere, measured in the actual scale of the spacetime. Without that decision the same rule gives a different theory, conformal gravity. Several attempts to derive the decision from the series' own structures failed; each turned out to be the same assumption in another form. So it is recorded as an assumption of its own. It explains why the cosmological constant is a fixed number, but for the same reason a measurement that dark energy never changes would not count as separate evidence for it. It would be refuted if dark energy were found to change, or if gravity were found to deviate from Einstein's equations where the rule applies.
\end{postulate}

\section{Dark energy and the frozen point}

These results identify dark energy with the length of one fixed direction in the completed algebra.

We start from a fact carried over from earlier in the series. For a gravitational field on flat spacetime, the vacuum equations amount to one tractor, the frozen point $I$, being parallel (covariantly constant), and with a cosmological constant $\Lambda$ its square length is \[ Q(I) = \Lambda/3 \]

The rotors that leave it fixed form the de Sitter group $SO(1,4)$, or the anti-de Sitter group when $\Lambda < 0$.

\begin{assembled}[Imported: the frozen point and its geometry]\label{c:dark-energy.imported-frozen-point-geometry}
This fact collects earlier results about the frozen point, a fixed direction in the six-dimensional conformal model attached to flat spacetime. Einstein's vacuum equations amount to this direction staying constant as the gravitational gauge fields carry it from place to place, and its length is set by the cosmological constant. The symmetries that leave it alone form the de Sitter group, or the anti-de Sitter group when the cosmological constant is negative. The fact also writes down the gauge field that says what constant means here.
\emph{First proved in} T10, Fact 2.1 (\href{https://doi.org/10.5281/zenodo.22722496}{doi:10.5281/zenodo.22722496}).
\end{assembled}
\noindent\emph{Rests on:} \ref{c:dark-energy.imported-frozen-point-geometry-external}, \ref{c:interior.covariant-form-vacuum-equations}, \ref{c:interior.imported-tractor-connection}, \ref{c:interior.non-null-stabilisers-section-they-act}, \ref{c:interior.scale-tractor-leaves-cone}, \ref{c:interior.sign-fixed-internally}, \ref{c:interior.stabiliser-rotates-parabolic-reductive}.

A single direction is not yet enough to write down gravity's action, so the next step builds a partner for it. From the frozen point and the canonical tractor $X$ of the conformal geometry comes a second vector, perpendicular to $I$: \[ V = X - \frac{3}{\Lambda}\, I \]

Its length is also set by $\Lambda$, and its rate of change from point to point is exactly the frame that links the six-dimensional tractor space to spacetime.

\begin{lemma}[The Stelle-West vector of the frozen point]\label{c:dark-energy.stelle-west-vector-frozen-point}
Starting from the frozen point and one canonical fixed direction of the six-dimensional conformal model, this lemma builds a vector perpendicular to the frozen point whose length is set by the cosmological constant. Its rate of change is exactly the frame that links the conformal model to the physical directions of the gravitational field on flat spacetime. It is the vector that the Stelle-West form of gravity needs, and the lemma records its basic algebra.
\emph{First proved in} T10, Lem 2.5 (\href{https://doi.org/10.5281/zenodo.22722496}{doi:10.5281/zenodo.22722496}).
\end{lemma}
\noindent\emph{Rests on:} \ref{c:dark-energy.imported-frozen-point-geometry}.

That partner vector is precisely what an older reformulation of gravity, due to MacDowell and Mansouri and to Stelle and West, requires. Parallelism of $I$ reduces the conformal $SO(2,4)$ structure to the de Sitter group, and the action becomes a square of the curvature of the $\Cl(2,4)$ gauge field, with no inverse metric anywhere.

So it makes sense wherever the connection stays regular, but only while the cosmological constant is nonzero.

\begin{theorem}[The frozen point makes the Stelle-West vector]\label{c:dark-energy.frozen-point-makes-stelle-west-vector}
The frozen point is a fixed direction in the larger conformal space whose length encodes the cosmological constant. On the region where Einstein's equations hold with a nonzero cosmological constant, this result shows that the frozen point, together with one canonical geometric object, produces exactly the extra vector that a known reformulation of gravity, due to MacDowell and Mansouri and to Stelle and West, requires. Gravity's action can then be written as the square of the curvature as seen from the point at infinity, with no inverse metric anywhere, so it is defined wherever the connection stays regular. At the black hole's wall, where the connection is singular, its value diverges but its field equations still reach the wall; and the curvature-squared form itself exists only while the cosmological constant is nonzero.
\emph{First proved in} T10, Thm 3.1 (\href{https://doi.org/10.5281/zenodo.22722496}{doi:10.5281/zenodo.22722496}).
\end{theorem}
\noindent\emph{Rests on:} \ref{c:dark-energy.imported-frozen-point-geometry}, \ref{c:dark-energy.imported-macdowell-mansouri-theorem-stelle-west}, \ref{c:dark-energy.imported-walls-sphere-theorem}, \ref{c:dark-energy.stelle-west-vector-frozen-point}.

An action built from two fixed-length vectors raises a worry: their wobbles might be extra fields that matter could feel. In the action the two vectors are held to the constraints \[ Q(I) = \Lambda/3 , \qquad \ip{V}{I} = 0 , \qquad Q(V) = -3/\Lambda \] and counting shows nine possible wobbles and nine broken symmetries, matched one to one.

So every wobble is pure gauge, and the graviton gains no extra polarisation.

\begin{proposition}[The point at infinity has no physical fluctuation: the $(I,V)$-sector is pure gauge]\label{c:dark-energy.point-infinity-has-no-physical-fluctuation}
Gravity's action here, written for gauge fields on flat spacetime, contains two fixed-length vectors tied to the point at infinity. One might worry that their wobbles are extra physical fields to which matter could couple. This result counts the possible wobbles and the broken symmetries, finds nine of each, and shows they match one to one, so every wobble can be undone by a symmetry. The graviton therefore carries no extra polarisation from this sector. This holds within the action's own variational principle, with both vectors treated as fields held to their fixed lengths; that the action so fixed is the known MacDowell-Mansouri one rests on the earlier theorem's assembly-level step.
\emph{First proved in} T10, Prop 5.4 (\href{https://doi.org/10.5281/zenodo.22722496}{doi:10.5281/zenodo.22722496}).
\end{proposition}
\noindent\emph{Rests on:} \ref{c:dark-energy.frozen-point-makes-stelle-west-vector}.

If the frozen point cannot wobble, its length cannot drift, and that is what dark energy would need to evolve. On the Einstein region the frozen point is parallel and the connection preserves lengths, so \[ \partial\_a Q(I) = 2\ip{\nabla\_a I}{I} = 0 \] and dark energy has pressure equal to minus its energy density, $w = -1$.

Going further needs a stated condition: that the frozen point can be defined in regions full of matter so that it stays parallel there, which the obvious extension does not do.

\begin{proposition}[$w = -1$ on the vacuum region; the registration, and its condition]\label{c:dark-energy.vacuum-region-registration-condition}
Where space is empty apart from the cosmological constant, the frozen point's length cannot change as the gravitational gauge fields carry it across flat spacetime, so dark energy behaves exactly like a cosmological constant, with pressure equal to minus its energy density. The paper goes further, on a stated condition: if dark energy is this length, and if the gravitational field is of the kind that keeps the frozen point fixed in regions full of matter, measured in the spacetime's own scale (an assumption the series now adopts; its own transport rule does not do this), then dark energy never evolves. With that assumption the cosmological constant is simply a fixed number, the standard cosmological model's assumption, so finding that dark energy never changes would not be separate evidence for it. That is a rigid, testable stance, and a confirmed measurement of evolving dark energy would refute it.
\emph{First proved in} T10, Prop 4.2 (\href{https://doi.org/10.5281/zenodo.22722496}{doi:10.5281/zenodo.22722496}).
\end{proposition}
\noindent\emph{Rests on:} \ref{c:dark-energy.frozen-point-makes-stelle-west-vector}, \ref{c:dark-energy.imported-frozen-point-geometry}, \ref{c:synthesis.field-space-premise-parallel-frozen-point}.

Granting that condition turns the result into a registered prediction. Dark energy cannot evolve: $w = -1$ exactly, at every epoch.

Any confirmed measurement of a different value at any point in cosmic history would rule the idea out.

\begin{prediction}[Dark energy cannot evolve: w = -1]\label{c:pred.dark-energy-w-minus-one}
The framework says dark energy is not a substance that can change over time but a fixed geometric quantity. That means its equation of state, the ratio of its pressure to its energy density, must be exactly minus one at all times. This rests on one stated assumption about how the underlying geometry extends into regions filled with matter, an assumption the series now adopts for its law of gravity. Under that assumption dark energy behaves exactly as in the standard cosmological model, so a measurement that it never changes would agree with the series without being separate evidence for the assumption. Any confirmed measurement of a different value at any point in cosmic history would rule the idea out.
\emph{Kill condition:} A confirmed detection of $w \neq -1$ at any epoch, with no parameter available to absorb it.
\end{prediction}

The prediction says the length is fixed; the remaining steps ask what fixes it. A conserved number in a physical theory is exactly one of three things: dynamical (D), a parameter of the laws (P), or initial data (I) that varies between solutions; being kinematic (K) is a separate question.

So ruling out the dynamical and kinematic options does not leave initial data as the only possibility.

\begin{proposition}[The four alternatives are exhaustive; the three are not]\label{c:dark-energy.four-alternatives-are-exhaustive-three-are}
A number in a physical theory can be fixed by the dynamics, can be a setting of the laws themselves, or can vary from one solution to another as initial data. These three options exclude one another and cover every case, while being kinematic is a separate question. So ruling out kinematic and dynamical does not leave initial data as the only possibility: the number might be a parameter of the laws.
\emph{First proved in} T10, Prop 5.7 (\href{https://doi.org/10.5281/zenodo.22722496}{doi:10.5281/zenodo.22722496}).
\end{proposition}

Applying that sorting to the frozen point's length gives a definite answer. Under the action exactly as written, $\Lambda$ is fixed before variation, and every solution has \[ H\^{}2 = Q(I) = \Lambda/3 \] on its de Sitter region, so the length cannot change within a solution yet differs between theories.

It is a parameter of the laws, not initial data. Making it initial data would need an extra equation of motion that the action does not provide.

\begin{proposition}[Under the action as written, the norm is a parameter of the laws]\label{c:dark-energy.under-action-written-norm-parameter-laws}
The cosmological constant fixes the length of the frozen point. Under the action exactly as written, this length cannot fluctuate within any one solution, yet it differs between theories with different cosmological constants. That makes it a setting of the laws rather than initial data chosen for each universe. Making it initial data would need an extra equation of motion that the action does not provide.
\emph{First proved in} T10, Prop 5.8 (\href{https://doi.org/10.5281/zenodo.22722496}{doi:10.5281/zenodo.22722496}).
\end{proposition}
\noindent\emph{Rests on:} \ref{c:dark-energy.four-alternatives-are-exhaustive-three-are}, \ref{c:dark-energy.frozen-point-makes-stelle-west-vector}, \ref{c:dark-energy.imported-frozen-point-geometry}, \ref{c:dark-energy.point-infinity-has-no-physical-fluctuation}, \ref{c:standard-model.trichotomy-what-kind-number}.

One might still hope the series' own machinery, the transport that repairs places where division fails, could pick the value. When the cosmological constant is zero the frozen point is null, and the transport completes that degenerate space to signature $(2,4)$, supplying a canonical null partner $\varphi$ of $I$.

Every vector orthogonal to the Lorentz frame is then $aI + b\varphi$ with norm $2ab$, and a frozen point of cosmological constant $\Lambda$ is one with \[ 2ab = \Lambda/3 \] So the transport fixes everything about the frozen point except its norm.

\begin{theorem}[The transport constrains the frozen point only through its norm]\label{c:odd-extensions.transport-constrains-frozen-point-only-through}
When the cosmological constant is zero the frozen point becomes a vector of zero length, a degenerate spot that the series' transport completes into a larger, well-behaved space. That completion comes with a canonical plane and a preferred pair of directions in it. Every possible value of the cosmological constant corresponds to a point in this plane, and the available symmetries can move that point anywhere except to a point of different length. So the transport pins down everything about the frozen point except its length, which is the cosmological constant, and gives no reason to prefer a positive value over a negative one.
\emph{First proved in} T14a, Thm 8.2 (\href{https://doi.org/10.5281/zenodo.23232104}{doi:10.5281/zenodo.23232104}).
\end{theorem}
\noindent\emph{Rests on:} \ref{c:odd-extensions.imported-minimal-completion-degenerate-algebra-rule}, \ref{c:odd-extensions.twistor-algebra}.

The last step states the consequence plainly. Within the series the norm $\Lambda/3$, sign and magnitude together, is data the transport receives rather than data it produces.

This does not mean nothing could fix it, only that the transport does not.

\begin{corollary}[The cosmological constant is undetermined within the series]\label{c:odd-extensions.cosmological-constant-undetermined-within-series}
The series has a procedure, the transport, for dealing with places where its algebra breaks down. Running it at the point where the cosmological constant would be zero produces a definite structure, but it leaves the value of the cosmological constant, including whether it is positive or negative, entirely open. So within the series the cosmological constant is an input rather than an output. This does not mean nothing could fix it, only that the transport does not.
\emph{First proved in} T14a, Cor 8.3 (\href{https://doi.org/10.5281/zenodo.23232104}{doi:10.5281/zenodo.23232104}).
\end{corollary}
\noindent\emph{Rests on:} \ref{c:odd-extensions.transport-constrains-frozen-point-only-through}.

\section{Results used from elsewhere}
The proofs above use the following results of the series, stated briefly; each is proved in the paper named.

\begin{fact}[Imported: the frozen point and its geometry (external ingredients)]\label{c:dark-energy.imported-frozen-point-geometry-external}
The frozen-point fact takes its tractor connection and almost-Einstein scale tractor from Gover's work on almost-Einstein manifolds.
\emph{Source:} A. R. Gover, Almost Einstein and Poincar'e-Einstein manifolds in Riemannian signature, J. Geom. Phys. 60 (2010), 182--204..
\end{fact}

\begin{fact}[Imported: the MacDowell-Mansouri theorem, Stelle-West form]\label{c:dark-energy.imported-macdowell-mansouri-theorem-stelle-west}
The MacDowell-Mansouri functional of a de Sitter gauge field $A(a) = \Omega(a) - \hb^{-1}(a)\wedge V/\ell^2$ on flat spacetime, with a Stelle-West vector of fixed norm, splits into Gauss-Bonnet, Einstein-Cartan and cosmological terms, and its field equations are the Einstein-Cartan equations with cosmological constant $\Lambda$.
\emph{Source:} S. W. MacDowell and F. Mansouri, Unified geometric theory of gravity and supergravity, Phys. Rev. Lett. 38 (1977), 739--742.; K. S. Stelle and P. C. West, Spontaneously broken de Sitter symmetry and the gravitational holonomy group, Phys. Rev. D 21 (1980), 1466--1488.; D. K. Wise, MacDowell-Mansouri gravity and Cartan geometry, Class. Quantum Grav. 27 (2010), 155010..
\end{fact}

\begin{fact}[Imported: walls, and the sphere theorem]\label{c:dark-energy.imported-walls-sphere-theorem}
At a tame corank-one wall the Clifford sheaf continues, the metric has a simple pole with $\mathcal{L}^{\otimes 2} \cong N^*\Sigma$, the completion adjoins a Witt partner of the radical covector, and, under the Indifference principle, the transverse sphere is carried by boundary-mode bilinears.
\emph{Source:} T10, Fact 2.4 (\href{https://doi.org/10.5281/zenodo.22722496}{doi:10.5281/zenodo.22722496}).
\end{fact}

\begin{fact}[Covariant form of the vacuum equations]\label{c:interior.covariant-form-vacuum-equations}
On the invertible region of a gravitational field on flat spacetime, the vacuum field equations hold exactly when the null scale tractor $\nInf$ of the derived metric is parallel for the tractor gauge field $\Omega_{T}$.
\emph{Source:} T6, Thm 4.5 (\href{https://doi.org/10.5281/zenodo.22721990}{doi:10.5281/zenodo.22721990}).
\end{fact}

\begin{fact}[Imported: the tractor connection]\label{c:interior.imported-tractor-connection}
For a gravitational field on flat spacetime, the tractor bundle of its conformal class (the $\Cl(2,4)$ vectors $\sigma\nInf + \mu + \rho\nO$) carries a canonical metric-preserving connection, given in a chosen scale by an explicit formula involving the Schouten tensor.
\emph{Source:} T. N. Bailey, M. G. Eastwood and A. R. Gover, Thomas's structure bundle for conformal, projective and related structures, Rocky Mountain J. Math. 24 (1994), 1191--1217.; S. N. Curry and A. R. Gover, An introduction to conformal geometry and tractor calculus, with a view towards analysis on manifolds, in Asymptotic Analysis in General Relativity, London Mathematical Society Lecture Note Series 443, Cambridge University Press (2018), 86--170., \S 5.
\end{fact}

\begin{fact}[Non-null stabilisers, and the section they act on]\label{c:interior.non-null-stabilisers-section-they-act}
A non-null vector $I$ has the reductive ten-dimensional stabiliser $\mathfrak{so}(I^\perp)$, and its hyperplane section of the null quadric is de Sitter or anti-de Sitter space according to the sign of $Q(I)$.
\emph{Source:} T6, Prop 4.4 (\href{https://doi.org/10.5281/zenodo.22721990}{doi:10.5281/zenodo.22721990}).
\end{fact}

\begin{fact}[$\Lambda \neq 0$: the scale tractor leaves the cone]\label{c:interior.scale-tractor-leaves-cone}
For a gravitational field whose derived metric is Einstein with constant $\Lambda$, the section $\nInf + \tfrac{\Lambda}{6}\nO$ is parallel, every parallel section $(1, 0, \rho)$ arises this way, and its squared norm is $\Lambda/3$.
\emph{Source:} T6, Thm 4.7 (\href{https://doi.org/10.5281/zenodo.22721990}{doi:10.5281/zenodo.22721990}).
\end{fact}

\begin{fact}[The sign is fixed internally]\label{c:interior.sign-fixed-internally}
Within the series' conventions, the sign of the Schouten term of the tractor gauge field is forced by requiring that $\Lambda > 0$ give the de Sitter section.
\emph{Source:} T6, Lem 4.8 (\href{https://doi.org/10.5281/zenodo.22721990}{doi:10.5281/zenodo.22721990}).
\end{fact}

\begin{fact}[The stabiliser rotates from parabolic to reductive]\label{c:interior.stabiliser-rotates-parabolic-reductive}
The stabiliser of the frozen point is the Poincar\'e group for $\Lambda = 0$, the de Sitter group $SO(1,4)$ for $\Lambda > 0$ and the anti-de Sitter group $SO(2,3)$ for $\Lambda < 0$.
\emph{Source:} T6, Prop 4.10 (\href{https://doi.org/10.5281/zenodo.22721990}{doi:10.5281/zenodo.22721990}).
\end{fact}

\begin{fact}[Imported: the minimal completion of a degenerate algebra, and Rule 1]\label{c:odd-extensions.imported-minimal-completion-degenerate-algebra-rule}
Minimal nondegenerate completions of a space with radical of dimension $r$ add $r$ hyperbolic planes and are unique up to an isometry fixing the space; Rule 1 forces the increment $(+1,+1)$ for type N, while for type C the completed algebra forgets the sign of the adjoined generator.
\emph{Source:} T14a, Fact 8.1 (\href{https://doi.org/10.5281/zenodo.23232104}{doi:10.5281/zenodo.23232104}).
\end{fact}

\begin{fact}[The even conformal algebra is $\Cl(2,3)$]\label{c:odd-extensions.twistor-algebra}
The even part of $\Cl(2,4)$ is $\Cl(2,3) \cong \Cl(4,1) \cong \Mat(4, \C)$ and contains both $\Cl(1,3)$ and $\Cl(3,1)$, while signature reversal does not identify the odd extensions.
\emph{Source:} T14a, Thm 4.1 (\href{https://doi.org/10.5281/zenodo.23232104}{doi:10.5281/zenodo.23232104}).
\end{fact}

\begin{fact}[The trichotomy: what kind of number $N$ is]\label{c:standard-model.trichotomy-what-kind-number}
The number $N = M_{P}/H$ rescales while every derived discrete invariant stays fixed (proved), which is read at observation level as $N$ not being kinematic; it is imported as not dynamical, and under the action as written it is a parameter of the laws, becoming initial data only if the norm is promoted to a field.
\emph{Source:} T12, Prop 7.6 (\href{https://doi.org/10.5281/zenodo.22722669}{doi:10.5281/zenodo.22722669}).
\end{fact}

\section*{Conclusion}
The thread runs from the frozen point, through a polynomial action whose extra vectors carry no physics, to a prediction that dark energy never evolves.

No postulate is needed on this stretch, but the prediction rests on one stated condition, that the frozen point extends unchanged into regions full of matter. Under the action as written, the frozen point's length is a setting of the laws rather than initial data chosen for each universe.

\end{document}